\documentclass[aspectratio=169]{beamer}
\input{header}
\coursecode{JKSSB}

\title{MS Word High-Yield}
\subtitle{Lesson 23 --- Shortcuts, Review Tools and Mail Merge}
\author{JKSSB Computer Awareness}
\date{\today}

\begin{document}
\frame{\titlepage}

\begin{frame}{Why These Word Tools Matter}
  \begin{itemize}
    \item Exam items on MS Word repeat a small set of \key{review and editing
      tools}.
    \item The high-yield set: Find/Replace, spell check, hyperlinks, styles,
      comments and Track Changes.
    \item Also tested: \key{Page Break vs Section Break}, mail merge, and
      save/print/PDF export.
    \item Most items turn on an \key{exact shortcut} or a
      \key{near-neighbour distinction}.
  \end{itemize}
  \begin{block}{The one idea}
    Learn each tool as a shortcut plus the one thing it is \emph{not}.
  \end{block}
\end{frame}

\begin{frame}{Find and Replace}
  \begin{itemize}
    \item \key{Ctrl+F} opens the Find box to locate text in the document.
    \item \key{Ctrl+H} opens \key{Find and Replace}.
    \item Replace can change one match or use \key{Replace All} for every
      occurrence.
    \item Find also locates formatting, special characters and styles.
  \end{itemize}
  \begin{block}{Keep the pair straight}
    Ctrl+F only \emph{finds}; Ctrl+H \emph{finds and replaces}.
  \end{block}
\end{frame}

\begin{frame}{Spelling and Grammar}
  \begin{itemize}
    \item \key{F7} starts the spelling and grammar check.
    \item Word flags a misspelling with a red squiggle and grammar issues with
      other underlines.
    \item The \key{Editor} pane groups Spelling, Grammar and Refinements.
    \item ``Ignore'' skips a word once; ``Add to dictionary'' accepts it
      permanently.
  \end{itemize}
  \begin{alertblock}{Shortcut alert}
    F7 is spelling and grammar in Word. \muted{F5 and F9 belong to PowerPoint
    slide shows, not Word proofing.}
  \end{alertblock}
\end{frame}

\begin{frame}{Hyperlinks and Bookmarks}
  \begin{itemize}
    \item \key{Ctrl+K} inserts a \key{hyperlink}.
    \item A hyperlink can point to a web page, a file, or a place in the same
      document.
    \item A \key{bookmark} gives a location or selection a name.
    \item A hyperlink can jump to a bookmark; a bookmark alone is
      \muted{not visible in print}.
  \end{itemize}
  \begin{block}{Key difference}
    A \key{hyperlink} is the clickable link; a \key{bookmark} is the named
    target it points to.
  \end{block}
\end{frame}

\begin{frame}{Styles}
  \begin{itemize}
    \item A \key{style} is a named set of formatting: font, size, colour,
      spacing and indent.
    \item Applying a \key{Heading} style gives the paragraph an outline level.
    \item Styles let the \key{Navigation Pane} and a \key{Table of Contents}
      update automatically.
    \item Editing a style updates \key{every} paragraph that uses it.
  \end{itemize}
  \begin{block}{Definition}
    A style is reusable formatting stored under a name.
  \end{block}
\end{frame}

\begin{frame}{Comments}
  \begin{itemize}
    \item A \key{comment} is a note attached to a point in the document.
    \item Comments \key{do not change} the document text.
    \item Reviewers can reply to a comment and mark it resolved.
    \item Comments are metadata and appear in the margin or the Reviewing
      Pane.
  \end{itemize}
\end{frame}

\begin{frame}{Track Changes}
  \begin{itemize}
    \item \key{Track Changes} records insertions, deletions and formatting
      edits.
    \item Each change stores the \key{author's name} and the time --- revision
      metadata.
    \item The reviewer can \key{Accept} or \key{Reject} each change.
    \item Accepting a change keeps it permanently and removes its revision
      markup.
  \end{itemize}
  \begin{alertblock}{Trap alert}
    Accepting a change permanently removes the author's metadata for that
    change. (PYQ-12)
  \end{alertblock}
\end{frame}

\begin{frame}{Page Break vs Section Break}
  \begin{center}
    \begin{tabular}{p{0.20\textwidth}p{0.30\textwidth}p{0.36\textwidth}}
      \toprule
      \textbf{Point} & \textbf{Page Break} & \textbf{Section Break} \\
      \midrule
      What it does & Moves text to a new page & Starts a new section \\
      Formatting & Same section layout & New layout and formatting \\
      Headers/footers & Cannot differ & Can differ per section \\
      Insert & Ctrl+Enter (page break) & Layout $\rightarrow$ Breaks \\
      \bottomrule
    \end{tabular}
  \end{center}
  \vspace{0.2cm}
  \muted{Page Break changes the page; Section Break changes the section.}
\end{frame}

\begin{frame}{Section Break and Headers/Footers}
  \begin{itemize}
    \item A \key{Section Break} divides the document into independent
      \key{sections}.
    \item Each section can have its own \key{headers and footers}.
    \item Each section can also set its own page numbering, margins and
      orientation.
    \item A \key{Page Break} does not allow different headers or footers.
  \end{itemize}
  \begin{alertblock}{Trap alert}
    Only a \key{Section Break} --- never a Page Break --- enables different
    headers and footers in one document. (TRAP-13 / PYQ-13)
  \end{alertblock}
\end{frame}

\begin{frame}{Mail Merge: The Idea}
  \begin{itemize}
    \item \key{Mail Merge} combines a \key{main document} with a
      \key{data source}.
    \item Main document: the letter or template, with \key{merge fields} as
      placeholders.
    \item Data source: a list of records such as names and addresses.
    \item Result: one \key{personalised copy} per record.
  \end{itemize}
  \begin{block}{Definition}
    Mail Merge combines a standard letter with a database of addresses.
    (PYQ-37, PYQ-47)
  \end{block}
\end{frame}

\begin{frame}{Mail Merge: Steps}
  \begin{exampleblock}{Typical mail-merge steps}
    \begin{enumerate}
      \item Create or select the main document (the letter).
      \item Select recipients (the data source).
      \item Insert merge fields such as name and address.
      \item Preview the merged letters.
      \item Finish and Merge --- print or send by email.
    \end{enumerate}
  \end{exampleblock}
\end{frame}

\begin{frame}{Saving, Printing and PDF Export}
  \begin{itemize}
    \item \key{Ctrl+S} saves the document, normally as a \key{DOCX} file.
    \item \key{Ctrl+P} opens the print options.
    \item \key{Save As} can change the file name, location and file format.
    \item \key{Export} or \key{Save As PDF} creates a
      \key{PDF} --- Portable Document Format --- file.
  \end{itemize}
  \begin{block}{PDF}
    Portable Document Format: a \key{fixed-layout} file that looks the same on
    every device, good for sharing and printing.
  \end{block}
\end{frame}

\begin{frame}{Shortcut Near-Neighbours}
  \begin{center}
    \begin{tabular}{ll}
      \toprule
      \textbf{Action} & \textbf{Shortcut} \\
      \midrule
      Find & Ctrl+F \\
      Find and Replace & Ctrl+H \\
      Spelling and grammar & F7 \\
      Insert a hyperlink & Ctrl+K \\
      Save & Ctrl+S \\
      Print & Ctrl+P \\
      Remove document-window split & Alt+Shift+C \\
      \bottomrule
    \end{tabular}
  \end{center}
\end{frame}

\begin{frame}{Trap Alert: Word Facts to Keep Straight}
  \begin{itemize}
    \item A \key{Section Break}, not a Page Break, enables different
      headers/footers. (TRAP-13)
    \item \key{Track Changes} records revision metadata; accepting a change
      removes that change's metadata. (PYQ-12)
    \item A \key{bookmark} is a named target; a \key{hyperlink} (Ctrl+K) is the
      clickable link.
    \item \key{F7} is spelling/grammar; F5 and F9 are PowerPoint slide-show
      keys. (TRAP-32)
    \item \key{Mail Merge} combines a letter with an address database.
      (TRAP-29)
  \end{itemize}
\end{frame}

\begin{frame}{Lesson 23: Recall}
  \begin{itemize}
    \item Find: \key{Ctrl+F}; Find and Replace: \key{Ctrl+H}; Replace All
      changes every match.
    \item Spell check: \key{F7}; hyperlink: \key{Ctrl+K}; save: Ctrl+S; print:
      Ctrl+P.
    \item \key{Styles} are named formatting; Heading styles drive the
      Navigation Pane and a Table of Contents.
    \item \key{Comments} do not change text; \key{Track Changes} records edits
      with author metadata.
    \item \key{Section Break} --- not Page Break --- enables different
      headers/footers.
    \item \key{Mail Merge} combines a standard letter with an address database;
      \key{Save As PDF} exports a fixed-layout file.
  \end{itemize}
\end{frame}
\end{document}
